%=====================================================================
%  Syllabus-Grounded Contextual Decoding for Code-Switched
%  Classroom Speech Transcription and Automated Note Generation
%
%  IEEEtran conference format. Target length: 6 pages.
%  COMPILE:  pdflatex ClassScribe_Research_Paper.tex   (run twice)
%
%  TO INSERT ARTWORK: replace \figplace{...}{...} with
%      \includegraphics[width=\columnwidth]{figures/yourfile.png}
%  full width:  \includegraphics[width=\textwidth]{figures/yourfile.png}
%=====================================================================
\documentclass[conference]{IEEEtran}
\IEEEoverridecommandlockouts

\usepackage{cite}
\usepackage{amsmath,amssymb}
\usepackage{array}
\usepackage{graphicx}

% centred fixed-width column, used for the author grid
\newcolumntype{C}[1]{>{\centering\arraybackslash}p{#1}}
\usepackage{textcomp}
\usepackage{xcolor}
\usepackage{url}

% ---- placeholder boxes that RESERVE SPACE for artwork ----------------
\newcommand{\figplace}[2]{%
  \setlength{\fboxsep}{0pt}%
  \fbox{\begin{minipage}[c][#1][c]{\dimexpr\columnwidth-2\fboxrule\relax}
    \centering\scriptsize\itshape\color{gray}#2
  \end{minipage}}%
}
\newcommand{\figplacewide}[2]{%
  \setlength{\fboxsep}{0pt}%
  \fbox{\begin{minipage}[c][#1][c]{\dimexpr\textwidth-2\fboxrule\relax}
    \centering\scriptsize\itshape\color{gray}#2
  \end{minipage}}%
}
\newcommand{\figpanel}[3]{%
  \setlength{\fboxsep}{0pt}%
  \fbox{\begin{minipage}[c][#2][c]{#1}
    \centering\scriptsize\itshape\color{gray}#3
  \end{minipage}}%
}

\begin{document}

\title{Syllabus-Grounded Contextual Decoding for\\
Code-Switched Classroom Speech Transcription and\\
Automated Study-Note Generation}

%--------------------------------------------------------------------
% AUTHORS -- edit as required.
%--------------------------------------------------------------------
% Six authors laid out as a balanced 3 x 2 grid. Using an explicit tabular
% instead of \and, because \and wraps automatically and strands the last row.
\author{%
\small
\begin{tabular}{C{0.31\textwidth}C{0.31\textwidth}C{0.31\textwidth}}
Meet Patel\newline
Information Technology\newline
SVKM's DJSCE\newline
Mumbai, India\newline
meet2005pokar@gmail.com
&
Vinay Vora\newline
Information Technology\newline
SVKM's DJSCE\newline
Mumbai, India\newline
voravinay6@gmail.com
&
Manav Chudasama\newline
Information Technology\newline
SVKM's DJSCE\newline
Mumbai, India\newline
manavchudasama2005@gmail.com
% \tabularnewline (not \\) ends the ROW: inside a p-column, \\ is only an
% in-cell line break, which is why the optional spacing was being ignored.
\tabularnewline[3ex]
Pranay Mistry\newline
Information Technology\newline
SVKM's DJSCE\newline
Mumbai, India\newline
pranaymistry21@gmail.com
&
Mohit Nippanikar\newline
Information Technology\newline
SVKM's DJSCE\newline
Mumbai, India\newline
rnippanikar7@gmail.com
&
Neha Katre\newline
Information Technology\newline
SVKM's DJSCE\newline
Mumbai, India\newline
neha.katre@djsce.ac.in
\end{tabular}%
}

\maketitle

%=====================================================================
\begin{abstract}
Classroom instruction in Indian higher education is routinely delivered in
code-switched Hindi--English, with technical terminology in English embedded in a
Hindi matrix. General-purpose speech recognition transcribes such lectures
poorly: domain terms are misrecognised and English vocabulary is frequently
rendered in Devanagari rather than Latin script, making transcripts unusable as
study material. This paper presents ClassScribe, a system that records lectures,
transcribes them and generates structured study notes with a subject-wise
knowledge assistant. Its technical contribution is Syllabus-Grounded Contextual
Decoding (SGCD), a training-free method that conditions a pre-trained decoder on
course-syllabus material rendered as fluent code-mixed narration rather than as
terminology lists, selected per utterance by sparse retrieval. On the MUCS 2021
Hindi--English corpus, narrative conditioning reduces word error rate by 6.23
points and technical-term error by 55.5\% relative on lecture-length spans,
whereas terminology-list conditioning degrades word error rate by 19.39 points
despite improving term recognition. A mismatched-syllabus control shows most of
the gain arises from orthographic and register conditioning rather than syllabus
content, a distinction not previously isolated for code-switched lecture ASR.
\end{abstract}

\begin{IEEEkeywords}
Automatic Speech Recognition, Code-Switching, Contextual Decoding, Hindi--English,
Lecture Transcription, Educational Technology, Note Generation.
\end{IEEEkeywords}

%=====================================================================
\section{Introduction}

Students in Indian higher education rely heavily on lecture notes for revision,
yet manual note-taking forces a trade-off: a student who writes cannot fully
attend, and one who attends retains no durable record. Automatic transcription
followed by automatic structuring is an attractive alternative, and has become
plausible with large-scale weakly supervised speech models such as Whisper
\cite{radford2023whisper}.

The Indian classroom makes this harder than the English lecture-transcription
literature suggests. Instruction is delivered in a code-switched register whose
matrix language is Hindi while the embedded technical vocabulary is English. A
sentence such as ``\textit{is tutorial mein hum file permissions samjhenge}''
carries its pedagogical content entirely in the English fragment. Two failure
modes follow. First, technical terminology---precisely what a student needs---is
the vocabulary least well covered by general-purpose acoustic models. Second, the
model must choose an orthography for each token, and a multilingual model
inclined toward Hindi transliterates English terms into Devanagari, producing
forms such as \textit{skreen} for \textit{screen}. Such a transcript is not
merely penalised by metrics; it is materially less useful, since technical terms
are no longer searchable in canonical form.

A resource addressing both failure modes already exists in every institution and
is systematically unused: the course syllabus. It enumerates, before delivery,
the topics to be covered and the vocabulary in which they are expressed. It needs
no annotation and is written by the instructor. This paper asks whether that
document can be supplied to a pre-trained decoder as auxiliary context, and what
form it must take to help rather than harm.

The question is not rhetorical. Prior work shows that naively injecting target
terms into the decoder's context channel improves those terms while degrading
everything else, worsening aggregate error on most datasets evaluated
\cite{sun2023contextual}. The implicated mechanism is distribution mismatch: the
channel is trained to receive the transcript of preceding speech, not a keyword
enumeration. Any method supplying syllabus content must respect that form.

Our contributions are: (i) \textbf{ClassScribe}, a system that records lectures,
produces code-switched transcripts, derives study notes and learning outcomes,
and exposes a conversational interface over a student's own lecture corpus;
(ii) \textbf{SGCD}, a training-free decoding procedure rendering syllabus units as
code-mixed narration matching the target dual-script convention, selected by
character $n$-gram TF--IDF retrieval over a first-pass hypothesis and protected by
a stability safeguard; (iii) a controlled evaluation separating the contribution
of \textit{representation format} from \textit{semantic content} using a
mismatched-syllabus control; and (iv) characterisation of the interaction between
context utility and utterance duration, with cross-model validation.

%=====================================================================
\section{Literature Survey}

\subsection{Lecture Capture and Code-Switched ASR}

Commercial services such as Otter.ai and platform-native live captions provide
real-time transcription and have been adopted in higher education, but are
optimised for monolingual conversational speech, offer limited configurability for
domain vocabulary, and process audio on remote infrastructure, raising
data-governance concerns when classroom recordings including student
contributions leave institutional control.

Code-switching, the alternation between languages within an utterance, is
pervasive in Indian educational discourse and is documented as difficult for
systems trained on monolingual corpora \cite{sitaram2019survey}. The MUCS 2021
challenge \cite{diwan2021mucs} targeted low-resource Indian-language ASR and
included a code-switching subtask with Hindi--English data drawn from spoken
technical tutorials, released as OpenSLR SLR104 \cite{openslr104}. To our
knowledge it is the only public, sentence-aligned Hindi--English code-switched
corpus whose audio is technical instructional content, making it the natural
vehicle for a syllabus-conditioned method.

An under-discussed property of code-switched Indic transcription is orthographic
ambiguity. SLR104 references follow a dual-script convention---Hindi in
Devanagari, English technical terms in Latin. A model given no indication of this
convention must guess, and its errors are systematic rather than random, hence in
principle correctable by conditioning. We quantify this with a script-fidelity
metric in Section~V.

\subsection{Contextual Adaptation and Research Gaps}

Contextual biasing steers recognition toward known relevant terms. Earlier
architectures used shallow fusion or attention over a biasing vocabulary, as in
Deep Context \cite{pundak2018deepcontext} and trie-based deep biasing for
streaming transducers \cite{le2021contextualized}. These require access to model
internals and usually training of a biasing module, placing them beyond
institutions deploying a pre-trained model without ML infrastructure.

For encoder--decoder models an alternative exists: conditioning the decoder on
preceding textual context at inference time, with no parameter updates. Peng et
al. \cite{peng2023prompting} showed this channel can elicit untrained task
behaviours. Sun et al. \cite{sun2023contextual} evaluated it for rare-word
biasing and reported the result motivating our design: a biasing list reduced
rare-word error from 23.7\% to 18.0\% and out-of-vocabulary error from 60\% to
37.1\%, but increased error on non-biased words on every dataset tested, with
aggregate error worsening on six of eleven.

Three gaps follow. Existing training-free methods are evaluated with term-list
representations known to degrade aggregate accuracy, and the representation
question is not itself treated as an object of study. Evaluations are almost
entirely monolingual English, leaving the code-switched case---where orthography
compounds vocabulary---unexamined. Finally, published evaluations rarely include
a control distinguishing a genuine content effect from a format effect; without
one, an improvement attributed to domain knowledge may be attributable to the
model being nudged into the correct output register. SGCD addresses all three.

%=====================================================================
\section{Proposed Methodology}

\subsection{System Overview}

ClassScribe begins with classroom audio and a course syllabus and terminates in
structured study material. The pipeline comprises audio acquisition and
segmentation; syllabus ingestion; first-pass transcription; contextual unit
retrieval; syllabus-grounded second-pass decoding with a stability safeguard; and
note synthesis with indexed retrieval for conversational query. Fig.~1 shows the
system architecture and the interfaces between these modules.

\begin{figure*}[htbp]
\centering
\figplacewide{6.2cm}{[ Fig. 1 --- System Architecture of ClassScribe ]\\[3pt]
Replace with: \texttt{\textbackslash includegraphics[width=\textbackslash textwidth]\{figures/fig1\_architecture.png\}}\\[4pt]
Suggested contents: recording client $\rightarrow$ segmentation $\rightarrow$ first-pass decoder $\rightarrow$
syllabus index (TF--IDF) $\rightarrow$ context assembly $\rightarrow$ second-pass decoder $\rightarrow$
stability safeguard $\rightarrow$ note synthesis / outcome extraction $\rightarrow$ vector store $\rightarrow$ chat interface}
\caption{System Architecture of ClassScribe}
\label{fig:arch}
\end{figure*}

Audio is captured at 16~kHz mono and partitioned into spans bounded by the
30-second encoder receptive field, so each span is decoded in a single forward
pass and the auxiliary context applies uniformly across it. Section~V-D shows
span duration is not an incidental detail but a determinant of whether
conditioning works at all; spans of approximately 25~seconds are used.

\subsection{Syllabus Representation}

Each course is represented as an ordered list of units, a unit $u_i$ being the
triple $(t_i, d_i, K_i)$ of title, narrative description and terminology set. The
design of $d_i$ is the core of the method, governed by three constraints derived
from the decoder's context behaviour.

\textit{Register and orthography.} $d_i$ is written as flowing instructional
narration in the target dual-script convention, Hindi function words in
Devanagari and English technical terms in Latin. This conveys the expected output
convention in addition to the expected vocabulary.

\textit{Terminal weighting.} Later context tokens influence generation more than
earlier ones. The most distinctive terminology is therefore placed at the end of
$d_i$, and multiple units are ordered by ascending relevance so the most relevant
terminates the context.

\textit{Budget.} The channel admits a bounded token count. Measured with the
model's own tokeniser, code-mixed narration costs approximately three tokens per
word, so each $d_i$ is authored at 60--105 tokens and the assembled context is
truncated from the left at 200 tokens, preserving the terminal region.

\subsection{Two-Pass Contextual Retrieval}

Choosing which units to supply requires knowing the utterance topic, which is
circular at decode time. An unconditioned first pass yields a hypothesis $h$ that
need not be accurate, only lexically indicative. Each unit is indexed as
$t_i \Vert d_i \Vert K_i$ under a character $n$-gram TF--IDF representation with
$n\in[3,5]$, and scored by cosine similarity:
\begin{equation}
s(u_i, h) = \frac{\phi(t_i \Vert d_i \Vert K_i)\cdot\phi(h)}
{\lVert \phi(t_i \Vert d_i \Vert K_i)\rVert\, \lVert \phi(h)\rVert}
\end{equation}
where $\phi(\cdot)$ is the TF--IDF projection. The $k$ highest-scoring units are
emitted in ascending order of $s$.

Character $n$-grams are preferred to word tokens for two reasons specific to this
setting: code-mixed text exhibits inconsistent transliteration, so word matching
fails on orthographic variants; and the first-pass hypothesis is by construction
noisy, requiring a representation that degrades gracefully. A sparse index is also
negligible against a neural encoder, needing no model download and operating in
milliseconds. Fig.~2 shows the resulting workflow.

\begin{figure}[htbp]
\centering
\figplace{5.0cm}{[ Fig. 2 --- Syllabus-Grounded Contextual Decoding Workflow ]\\[3pt]
Replace with:\\ \texttt{\textbackslash includegraphics[width=\textbackslash columnwidth]\{figures/fig2\_sgcd.png\}}}
\caption{Syllabus-Grounded Contextual Decoding Workflow}
\label{fig:sgcd}
\end{figure}

\subsection{Stability Safeguard}

Conditioning is not uniformly beneficial per utterance. Inspection shows two
opposing effects: it repairs degenerate repetition present in unconditioned
output, and it also induces such repetition, where generation continues the
register of the supplied narration instead of terminating. The system therefore
accepts the conditioned hypothesis only when decoding remained stable. With
$\ell_c,\ell_u$ the mean token log-probabilities of conditioned and unconditioned
hypotheses, $r_c$ the compression ratio and $|\cdot|$ length in words, the
unconditioned hypothesis is substituted when any of
\begin{equation}
\ell_c < \ell_u - \delta, \qquad
r_c > \rho, \qquad
|h_c| > \lambda\,|h_u|
\end{equation}
holds, with $\delta,\rho,\lambda$ fitted on a development partition and frozen
before evaluation.

\subsection{Note Synthesis and Conversational Retrieval}

The corrected transcript is segmented topically and converted into a structured
note comprising a topic hierarchy, definitions of encountered terminology and
learning outcomes expressed as competency statements. Because the syllabus unit
for each span is already known from retrieval, notes are attached to the unit
they instantiate, yielding coverage tracking: the system reports which units have
been delivered and which remain outstanding. Transcript spans and notes are
embedded in a vector index keyed by subject and lecture; the conversational
interface answers queries by retrieving relevant spans and grounding its response
in them, citing the originating lecture and timestamp so any statement can be
verified against the recording.

The complete method requires no training, annotation or auxiliary acoustic model,
and costs one supplementary decoding pass, since retrieval reuses the first-pass
hypothesis the baseline produces anyway.

%=====================================================================
\section{System Implementation}

The system is a client--server application with a Python backend. The
transcription service wraps a pre-trained multilingual encoder--decoder speech
model executed through a hardware-accelerated runtime for Apple silicon
\cite{hannun2023mlx}, enabling the full evaluation suite to run on a single
consumer laptop at 5.5--15.7$\times$ real time. Sparse retrieval uses
scikit-learn; scoring and alignment use \texttt{jiwer}. Persistence is provided by
a relational store for lecture metadata, syllabi and notes, with a vector index
for transcript spans. Because transcription executes locally, classroom audio need
not leave institutional infrastructure, addressing the governance objection to
cloud services noted in Section~II-A. Fig.~3 shows the dashboard, the recording
and transcription view, the generated notes, and the conversational interface.

\begin{figure*}[htbp]
\centering
\figpanel{0.48\textwidth}{4.6cm}{[ (a) Dashboard --- subject overview ]\\[2pt]
\texttt{\textbackslash includegraphics[width=0.48\textbackslash textwidth]\{figures/fig3a\_dashboard.png\}}}
\hfill
\figpanel{0.48\textwidth}{4.6cm}{[ (b) Lecture recording and transcript ]\\[2pt]
\texttt{\textbackslash includegraphics[width=0.48\textbackslash textwidth]\{figures/fig3b\_recording.png\}}}\\[6pt]
\figpanel{0.48\textwidth}{4.6cm}{[ (c) Generated notes and learning outcomes ]\\[2pt]
\texttt{\textbackslash includegraphics[width=0.48\textbackslash textwidth]\{figures/fig3c\_notes.png\}}}
\hfill
\figpanel{0.48\textwidth}{4.6cm}{[ (d) Subject-wise conversational interface ]\\[2pt]
\texttt{\textbackslash includegraphics[width=0.48\textbackslash textwidth]\{figures/fig3d\_chat.png\}}}
\caption{ClassScribe interface: (a) subject dashboard, (b) lecture recording and
code-switched transcript, (c) generated study notes and learning outcomes,
(d) subject-wise conversational assistant.}
\label{fig:ui}
\end{figure*}

%=====================================================================
\section{Experimental Evaluation}

\subsection{Dataset and Protocol}

Evaluation uses the Hindi--English code-switched test partition of MUCS 2021
subtask-2 \cite{diwan2021mucs}, released as OpenSLR SLR104 \cite{openslr104}
under CC~BY-SA~4.0. The audio consists of spoken technical tutorials, so a
syllabus is a natural rather than contrived context source. The partition
comprises 30 lectures; after filtering to utterances of 2--28~s with at least
four reference words, 3073 utterances remain, totalling 5.06~hours.

All experiments are zero-shot; no component is trained on this corpus and the
training partition is never used, so the train/test overlap reported for this
corpus does not apply. Lectures are partitioned disjointly: 30\% form a
development partition used solely for selecting decoding configuration, and the
remainder an evaluation partition decoded once, after configuration was frozen.
Stratified sampling proportional to lecture size with a fixed seed yields 60
development and 150 evaluation utterances.

Because the corpus encodes no topic identifiers, each lecture's subject was
identified from its opening utterances, which announce the tutorial topic; those
utterances are excluded from the evaluation sample, so no scored utterance
contributed to constructing any syllabus. Seven syllabi covering 62 units were
authored from lecture titles alone, with no access to any reference transcript of
any scored utterance.

\subsection{Evaluation Metrics}

Aggregate word error rate (WER) cannot distinguish improvement on domain
terminology from degradation elsewhere, so we decompose it. With $R$ the
reference tokens and $K \subset R$ those matching the syllabus terminology
vocabulary, \textbf{K-WER} is the error rate over $K$ and \textbf{U-WER} that
over $R \setminus K$. \textbf{Script fidelity} is the proportion of Latin-script
reference tokens recovered exactly, measuring adherence to the dual-script
convention. We also report \textbf{WER\textsubscript{sa}}, a
transliteration-tolerant variant reducing both scripts to a common consonant
skeleton, so a term recognised correctly but written in the other script is not
penalised; it is a lower bound on error where strict WER is an upper bound.
Significance uses paired bootstrap resampling over utterances
\cite{koehn2004significance} with 10{,}000 resamples, reporting 95\% confidence
intervals on the difference in corpus error rate.

\subsection{Conditioning Strategies}

Three strategies form the primary comparison: \textbf{A}, unconditioned baseline;
\textbf{B}, terminology enumeration, the syllabus term sets as a delimited list;
and \textbf{C}, narrative representation, identical content as fluent code-mixed
narration. Four further conditions serve as controls: \textbf{D}, retrieved
narrative (SGCD), restricted to the $k=3$ units retrieved per utterance;
\textbf{E}, mismatched syllabus, as D but drawn from a \textit{different} course,
identical in length, register and construction with only the content wrong;
\textbf{F}, oracle retrieval against the reference, an upper bound reported as
topline only; and \textbf{G}, a register control consisting of one generic
code-mixed sentence with no course content, isolating pure format conditioning.

\begin{table*}[htbp]
\caption{Comparison of Conditioning Strategies on Utterance-Level Segments (150 utterances, mean 5.7~s)}
\label{tab:main}
\centering
\begin{tabular}{|l|c|c|c|c|c|c|}
\hline
\textbf{Strategy} & \textbf{WER (\%)} & \textbf{WER\textsubscript{sa} (\%)} &
\textbf{K-WER (\%)} & \textbf{U-WER (\%)} & \textbf{Script Fid. (\%)} &
\textbf{$\Delta$WER vs. A [95\% CI]} \\
\hline
A --- Unconditioned            & 85.69  & 71.87  & 65.53 & 43.50 & 32.6 & --- \\
B --- Terminology enumeration  & 136.22 & 118.43 & \textbf{31.55} & 53.18 & 60.9 & $+50.53$ $[+16.55, +90.94]$ \\
C --- Narrative representation & 99.35  & 88.81  & 33.50 & \textbf{41.68} & 60.0 & $+13.66$ $[-7.70, +35.35]$ \\
\hline
G --- Register control (no content) & 102.32 & 90.48 & 55.83 & 44.38 & 41.6 & $+16.63$ $[-3.61, +36.66]$ \\
D --- Retrieved narrative (SGCD)    & 90.80  & 81.27 & 35.44 & 41.95 & 58.8 & $+5.11$ $[-13.24, +23.28]$ \\
E --- Mismatched syllabus           & 87.71  & 77.62 & 35.92 & 40.19 & 58.6 & $+2.02$ $[-13.98, +18.98]$ \\
F --- Oracle retrieval (topline)    & 95.61  & 85.64 & 29.61 & 40.53 & 62.1 & $+9.92$ $[-10.49, +30.23]$ \\
D + stability safeguard             & \textbf{71.50} & \textbf{61.28} & 35.44 & 41.14 & 58.1 & $\mathbf{-14.19}$ $[-26.55, -4.79]$ \\
\hline
\end{tabular}
\end{table*}

\subsection{Results}

Table~I reports the primary comparison on the 150-utterance sample, mean duration
5.7~s. Terminology enumeration reproduces the published failure mode
\cite{sun2023contextual} in a new language setting: K-WER falls from 65.53\% to
31.55\%, a 51.9\% relative improvement on exactly the targeted vocabulary, while
U-WER rises from 43.50\% to 53.18\% and aggregate WER degrades by 50.53 points
with a confidence interval excluding zero. Since non-terminology tokens dominate
any utterance, the aggregate outcome follows the degradation. Narrative
representation attains an equivalent terminology gain, K-WER 33.50\%, while
reversing the collateral damage: U-WER falls to 41.68\%, below baseline. The
direct contrast is $-36.88$ points $[-75.64, -5.18]$ favouring narrative form.
Representation, not content volume, governs whether syllabus context is usable.

The absolute error rates in Table~I are high. Two corpus properties account for
this: segment boundaries are quantised to whole seconds, so fragments of adjacent
utterances intrude; and the mean segment of 5.7~s affords a context mechanism
little material to influence. We therefore built a secondary set by concatenating
strictly adjacent segments of the same lecture into spans bounded by the encoder
receptive field, yielding 100 spans averaging 26.2~s and 55.1 reference words
against approximately 12 previously. This analysis followed the primary
evaluation and is reported as a post-hoc experiment.

\begin{table*}[htbp]
\caption{Comparison of Conditioning Strategies on Lecture-Length Spans (100 spans, mean 26.2~s)}
\label{tab:long}
\centering
\begin{tabular}{|l|c|c|c|c|c|c|}
\hline
\textbf{Strategy} & \textbf{WER (\%)} & \textbf{WER\textsubscript{sa} (\%)} &
\textbf{K-WER (\%)} & \textbf{U-WER (\%)} & \textbf{Script Fid. (\%)} &
\textbf{$\Delta$WER vs. A [95\% CI]} \\
\hline
A --- Unconditioned            & 43.46 & 28.05 & 48.04 & 34.64 & 48.0 & --- \\
B --- Terminology enumeration  & 62.86 & 47.69 & 23.00 & 46.42 & 70.3 & $+19.39$ $[+7.13, +33.06]$ \\
C --- Narrative representation & \textbf{37.23} & \textbf{24.17} & 21.38 & \textbf{32.42} & \textbf{71.6} & $\mathbf{-6.23}$ $[-10.59, -1.97]$ \\
\hline
D --- Retrieved narrative (SGCD) & 38.64 & 25.53 & \textbf{21.24} & 33.94 & 70.0 & $-4.83$ $[-9.21, -0.53]$ \\
E --- Mismatched syllabus        & 39.56 & 26.40 & 25.44 & 33.91 & 67.3 & $-3.91$ $[-8.53, +0.61]$ \\
\hline
\end{tabular}
\end{table*}

Table~II shows baseline error approximately halves, from 85.69\% to 43.46\%,
confirming the primary table's magnitudes were inflated by segmentation artefacts
rather than by the method. At this span length the picture is unambiguous:
narrative conditioning reduces WER by 6.23 points $[-10.59, -1.97]$, terminology
error falls 55.5\% relative, and script fidelity rises from 48.0\% to 71.6\%.
Terminology enumeration remains harmful at $+19.39$ points $[+7.13, +33.06]$, and
the contrast between representations, $-25.62$ $[-39.28, -13.50]$, is
significant. Every contrast in Table~II excludes zero. The method therefore
requires realistic lecture-length spans, a deployment constraint rather than a
limitation of validity, since a classroom system controls its own segmentation.

\subsection{Content Specificity}

The mismatched control decomposes the benefit into format and content components.
On lecture-length spans the matched syllabus yields $-4.83$ points and the
mismatched syllabus $-3.91$. Content specificity therefore accounts for $-0.92$
points $[-3.53, +1.34]$, approximately 19\% of the effect, with an interval
including zero. The register control G is decisive in interpretation: a single
generic code-mixed sentence with no course content raises script fidelity from
32.6\% to 41.6\%, and any syllabus-derived context, matched or not, raises it to
roughly 60\%. The dominant mechanism is communication of the target output
convention---English terminology in Latin script within a Devanagari
matrix---rather than of which specific terms will occur.

This is a negative result with respect to the motivating hypothesis and we report
it as such. Its practical implication is favourable: a deployment does not need an
accurate per-lecture syllabus to obtain most of the benefit, lowering the barrier
to adoption. Its scientific implication is that evaluations omitting a mismatched
control cannot distinguish semantic grounding from register conditioning.
Resolving a 0.92-point effect at the observed variance would require roughly 700
spans, the full five-hour partition, which we identify as the decisive follow-up.

\subsection{Cross-Model Validation and Safeguard}

Repeating the conditions on a smaller model of the same family preserves the
qualitative ordering (Table~III): enumeration is catastrophic, narrative
representation beneficial, and their relative ranking unchanged. The smaller model
is almost entirely script-blind, achieving 1.6\% script fidelity unconditioned and
misrecognising 99.5\% of terminology, and conditioning is the only intervention
recovering canonical orthography. At this scale the matched syllabus does
outperform the mismatched by 17.52 points, though the interval spans zero. A
plausible reading, advanced as hypothesis rather than finding, is that content
specificity matters when a model is too weak to produce terminology unaided,
whereas for a stronger model only the output convention remains to be supplied.

\begin{table}[htbp]
\caption{Cross-Model Validation (150 utterances)}
\label{tab:crossmodel}
\centering
\resizebox{\columnwidth}{!}{%
\begin{tabular}{|l|c|c|c|}
\hline
\textbf{Strategy} & \textbf{WER (\%)} & \textbf{K-WER (\%)} & \textbf{Script Fid. (\%)} \\
\hline
A --- Unconditioned            & 171.50 & 99.51 & 1.6 \\
B --- Terminology enumeration  & 330.76 & 40.29 & 46.7 \\
C --- Narrative representation & 148.81 & 80.58 & 20.5 \\
D --- Retrieved narrative      & \textbf{143.23} & \textbf{74.27} & \textbf{23.7} \\
E --- Mismatched syllabus      & 160.75 & 88.35 & 16.0 \\
\hline
\end{tabular}}
\end{table}

On utterance-level segments the safeguard reverted 12.7\% of hypotheses,
converting the $+5.11$ regression of condition D into a $-14.19$ improvement
$[-26.55, -4.79]$ and reducing degraded utterances from 34\% to 24\%. Two
limitations apply: the gain concentrates in a few catastrophically degenerate
utterances rather than spreading across the corpus, and the thresholds do not
transfer across scales---applied unchanged to the smaller model the safeguard
fired on 46\% of utterances and degraded results, so thresholds must be refitted
per checkpoint. On lecture-length spans the safeguard is unnecessary, as condition
C attains a significant improvement without it.

\subsection{Threats to Validity}

The evaluation uses one corpus of scripted tutorial narration, more fluent than
spontaneous lecturing. Absolute error rates are inflated by whole-second segment
quantisation and are not comparable with published figures for this corpus. The
samples of 150 utterances and 100 spans are small, reflected in interval widths,
and the content-specificity question in particular is underpowered. Syllabi are
generated from lecture titles rather than published course outlines; although no
scored utterance contributed to their construction, syllabi authored by a language
model constitute a confound that institutionally published outlines would remove.

%=====================================================================
\section{Conclusion and Future Scope}

This paper presented ClassScribe, a system converting code-switched classroom
lectures into structured study notes and an interactive knowledge base, together
with SGCD, a training-free method for supplying syllabus material to a
pre-trained speech decoder. The principal finding concerns representation:
supplying terminology as an enumeration improves the enumerated terms by roughly
half while degrading everything else, so aggregate error worsens substantially,
whereas identical content rendered as fluent code-mixed narration achieves the
same terminology gain with no collateral damage, reducing word error by 6.23
points and terminology error by 55.5\% relative on lecture-length spans. The
determining factor is whether auxiliary context matches the distribution the
decoder was trained to receive. Controlled analysis further shows most of this
benefit conveys the dual-script output convention rather than syllabus semantics,
a syllabus from an unrelated course capturing roughly 81\% of the improvement.

Future work should re-examine content specificity at the scale identified in
Section~V-E, evaluate on spontaneous rather than scripted classroom speech,
replace the safeguard with a threshold-free criterion avoiding per-checkpoint
calibration, and assess the note-generation and conversational components against
learning outcomes with student participants, since transcription accuracy is a
proxy for, not a measure of, educational utility.

%=====================================================================
\begin{thebibliography}{00}

\bibitem{radford2023whisper}
A. Radford, J. W. Kim, T. Xu, G. Brockman, C. McLeavey, and I. Sutskever,
``Robust Speech Recognition via Large-Scale Weak Supervision,'' in
\textit{Proc. 40th Int. Conf. on Machine Learning (ICML)}, 2023, pp. 28492--28518.

\bibitem{sun2023contextual}
G. Sun, X. Zheng, C. Zhang, and P. C. Woodland, ``Can Contextual Biasing Remain
Effective with Whisper and GPT-2?,'' in \textit{Proc. Interspeech}, 2023,
pp. 1289--1293.

\bibitem{peng2023prompting}
P. Peng, B. Yan, S. Watanabe, and D. Harwath, ``Prompting the Hidden Talent of
Web-Scale Speech Models for Zero-Shot Task Generalization,'' in
\textit{Proc. Interspeech}, 2023, pp. 396--400.

\bibitem{hannun2023mlx}
A. Hannun, J. Digani, A. Katharopoulos, and R. Collobert, ``MLX: Efficient and
Flexible Machine Learning on Apple Silicon,'' 2023. [Online]. Available:
\url{https://github.com/ml-explore/mlx}

\bibitem{diwan2021mucs}
A. Diwan, R. Vaideeswaran, S. Shah, A. Singh, S. Raghavan, S. Khare, V. Unni,
S. Vyas, A. Rajpuria, C. Yarra, A. Mittal, P. K. Ghosh, P. Jyothi, K. Bali,
V. Seshadri, S. Sitaram, S. Bharadwaj, J. Nanavati, R. Nanavati, and
K. Sankaranarayanan, ``MUCS 2021: Multilingual and Code-Switching ASR Challenges
for Low Resource Indian Languages,'' in \textit{Proc. Interspeech}, 2021,
pp. 2446--2450.

\bibitem{le2021contextualized}
D. Le, M. Jain, G. Keren, S. Kim, Y. Shi, J. Mahadeokar, J. Chan, Y. Shangguan,
C. Fuegen, O. Kalinli, Y. Saraf, and M. L. Seltzer, ``Contextualized Streaming
End-to-End Speech Recognition with Trie-Based Deep Biasing and Shallow Fusion,''
in \textit{Proc. Interspeech}, 2021, pp. 1772--1776.

\bibitem{openslr104}
``OpenSLR SLR104: Multilingual and Code-Switching ASR Challenge Dataset,
Subtask-2,'' Open Speech and Language Resources, 2021. [Online]. Available:
\url{https://www.openslr.org/104/}

\bibitem{sitaram2019survey}
S. Sitaram, K. R. Chandu, S. K. Rallabandi, and A. W. Black, ``A Survey of
Code-switched Speech and Language Processing,'' \textit{arXiv preprint
arXiv:1904.00784}, 2019.

\bibitem{pundak2018deepcontext}
G. Pundak, T. N. Sainath, R. Prabhavalkar, A. Kannan, and D. Zhao, ``Deep
Context: End-to-End Contextual Speech Recognition,'' in \textit{IEEE Spoken
Language Technology Workshop (SLT)}, 2018, pp. 418--425.

\end{thebibliography}

\end{document}
